\chapter{Your Python Toolkit and Environment Setup}
\label{ch:python}

\chapterguide%
  {Create folders and run basic \textbf{terminal commands}.}
  {organize data and notebooks, configure an environment, install dependencies, and run labs reproducibly.}

\section{The course project layout}

Use \texttt{UTP-MACHINE-LEARNING/} as the project root, with these folders:

\noindent\textbf{Datasets}
\begin{lstlisting}[style=mlterminal]
all_datasets/
  lab01-03_airports.csv
  lab01-03_routes.csv
  lab04-AirQuality.csv
  lab05-sensor_readings_24.csv
  lab06-ocr-handwritten-prescription.zip
  lab07-Tweets.csv
  lab08-insurance.csv
\end{lstlisting}

\noindent\textbf{Experiments, helper code, and generated evidence}
\begin{lstlisting}[style=mlterminal]
all_experiments/
  lab01_search_foundations_worked.ipynb
  lab02_informed_search_worked.ipynb
  lab03_search_engineering_worked.ipynb
  lab04_fuzzy_anfis_worked.ipynb
  lab05_robotics_sensor_control_worked.ipynb
  lab06_handwriting_ocr_worked.ipynb
  lab07_airline_sentiment_worked.ipynb
  lab08_genetic_algorithm_regression_worked.ipynb
  utils/
    core.py
    openflights.py
    plotting.py
    __main__.py
    test_search.py
    search_queries.csv
results/
  lab01/ ... lab08/
requirements.txt
\end{lstlisting}

\begin{keyidea}
The notebooks import reusable search code from \texttt{utils/}.
\texttt{results/} stores generated evidence and should not be edited by hand;
\texttt{requirements.txt} lists the packages installed in \texttt{.venv/}.
\end{keyidea}

{\footnotesize Keep the digits ZIP intact. Each lab has one worked notebook; \texttt{predict\_passenger.py} uses the saved Lab~2 model. Notebooks write models and cards to \texttt{models/} and evidence to \texttt{results/}; do not edit generated outputs.\par}

\section{Download the datasets}

Click a \textbf{dataset name} to download it, then extract the files into \textbf{\texttt{all\_datasets/}}.

\begin{mltable}
\footnotesize
\renewcommand{\arraystretch}{1.1}
\begin{tabularx}{\linewidth}{@{}p{0.40\linewidth}>{\centering\arraybackslash}p{0.09\linewidth}X@{}}
\toprule
\rowcolor{MLPaleBlue!92}
\textbf{Kaggle dataset} & \textbf{Lab} & \textbf{Use} \\
\midrule
\textbf{OpenFlights:}
\href{https://www.kaggle.com/datasets/open-flights/flight-route-database}{\textbf{Routes}}
\& \href{https://www.kaggle.com/datasets/open-flights/airports-train-stations-and-ferry-terminals/data}{\textbf{Airports}}
& 1--3 & Route search, locations and distances \\
\href{https://www.kaggle.com/datasets/fedesoriano/air-quality-data-set}
{\textbf{Air Quality}} & 4 & Fuzzy logic and ANFIS \\
\href{https://www.kaggle.com/datasets/uciml/wall-following-robot}
{\textbf{Wall-Following Robot}} & 5 & Sensor-based control \\
\href{https://www.kaggle.com/datasets/mrdude20/doctor-handwriting-recognition-dataset}
{\textbf{Doctor Handwriting}} & 6 & OCR evaluation \\
\href{https://www.kaggle.com/datasets/chandana890/twitter-airline-sentiment-dataset}
{\textbf{Twitter Airline Sentiment}} & 7 & Sentiment classification \\
\href{https://www.kaggle.com/datasets/mirichoi0218/insurance}
{\textbf{Medical Cost}} & 8 & GA regression \\
\bottomrule
\end{tabularx}
\end{mltable}

\subsection{Download steps}

\begin{mlsteps}
  \item Open the link, click \textbf{Download}, and \textbf{extract} the archive.
  \item Copy the required \textbf{CSV files or images} into \textbf{\texttt{all\_datasets/}}.
  \item Check that your notebook's \textbf{file paths} match the filenames, including \textbf{capitalization}.
\end{mlsteps}

\begin{pitfall}
For \textbf{\texttt{FileNotFoundError}}, check the working directory and filename. Fix the path rather than moving files into the project root.
\end{pitfall}

\section{Two ways to set up: plain Python or Anaconda}

\begin{keyidea}
Use one environment per project; activate it before installing or running packages.
\end{keyidea}

\section{Route A: install Python and configure PATH}

Download Python from the \href{https://www.python.org/downloads/}{\textbf{official downloads page}}.

\begin{description}[leftmargin=1.15in,style=nextline]
  \item[Windows] Run the installer, enable \textbf{\texttt{PATH}}, and reopen the terminal. Use \textbf{\texttt{py}}.
  \item[macOS] Run the \textbf{\texttt{.pkg} installer} and restart Terminal. Keep the \textbf{system Python} unchanged.
  \item[Linux] Install \textbf{Python 3} with your package manager. Do not replace the \textbf{system Python}.
\end{description}

\textbf{Windows PATH fix.} If \texttt{py} is missing, open \textbf{Edit environment variables for your account}, then \textbf{Path > Edit}. Add the installer's command directory, for example:

\begin{lstlisting}[style=mlterminal]
%LocalAppData%\Python\bin
\end{lstlisting}

\textbf{macOS PATH fix.} If the installer placed Python in \texttt{/usr/local/bin} but your terminal cannot find it, add that directory to your Z shell's \textbf{PATH}:

\begin{lstlisting}[style=mlterminal]
echo 'export PATH="/usr/local/bin:$PATH"' >> ~/.zshrc
source ~/.zshrc
\end{lstlisting}

For \textbf{Debian or Ubuntu}:

\begin{lstlisting}[style=mlterminal]
sudo apt update
sudo apt install python3 python3-venv python3-pip
\end{lstlisting}

\textbf{Verify} the installation and executable location:

\begin{lstlisting}[style=mlterminal]
# Windows
py --version
py -m pip --version
where py

# macOS or Linux
python3 --version
python3 -m pip --version
which -a python3
\end{lstlisting}

\section{Route A: create an isolated virtual environment}

\textbf{\texttt{venv}} keeps this project's packages separate from other projects. See the \href{https://docs.python.org/3/library/venv.html}{\textbf{venv documentation}}.

The root-level \textbf{\texttt{requirements.txt}} lists dependencies, including PyTorch (\texttt{torch}) for Labs~7--8. A CPU is sufficient.

\begin{lstlisting}[style=mlterminal]
numpy
pandas
scipy
scikit-learn
matplotlib
jupyter
nbformat
nbclient
nbconvert
pillow
pytesseract
pytest
\end{lstlisting}

\textbf{Create} the environment in the \textbf{project root}:

\begin{lstlisting}[style=mlterminal]
# Windows
py -m venv .venv

# macOS or Linux
python3 -m venv .venv
\end{lstlisting}

\textbf{Activate} it from the same folder:

\begin{lstlisting}[style=mlterminal]
# Windows (PowerShell)
.venv\Scripts\Activate.ps1

# macOS/Linux
source .venv/bin/activate
\end{lstlisting}

\textbf{Install dependencies} after activation:

\begin{lstlisting}[style=mlterminal]
# Windows
py -m pip install --upgrade pip
py -m pip install -r requirements.txt

# macOS or Linux
python3 -m pip install --upgrade pip
python3 -m pip install -r requirements.txt
\end{lstlisting}

\textbf{Install once; activate every session.} Reuse \textbf{\texttt{.venv}} without reinstalling unchanged requirements.

\begin{mlfigure}
\begin{tikzpicture}[
  font=\small,
  pystep/.style={draw=MLBlue, fill=white, rounded corners=2pt, align=center,
               inner sep=4pt, minimum width=3.3cm, minimum height=1.05cm,
               font=\scriptsize\ttfamily\color{MLNavy}},
  pyonce/.style={pystep, fill=MLPaleBlue},
  pyevery/.style={pystep, fill=MLPaleTeal, draw=MLTeal},
  pycap/.style={font=\scriptsize\bfseries\color{MLMuted}},
  pynote/.style={font=\scriptsize\color{MLMuted}, align=center},
  pygo/.style={-{Stealth[length=2.2mm]}, thick, draw=MLTeal}]

\node[pyonce]  (venv) at (0,0)     {py -m venv .venv};
\node[pyevery] (act)  at (4.05,0)  {activate};
\node[pyonce]  (inst) at (8.10,0)  {pip install -r\\requirements.txt};
\node[pyevery] (work) at (12.15,0) {run notebooks in\\all\_experiments/};

\draw[pygo] (venv) -- (act);
\draw[pygo] (act)  -- (inst);
\draw[pygo] (inst) -- (work);

\node[pycap] at (0,0.92)     {once per project};
\node[pycap] at (4.05,0.92)  {every session};
\node[pycap] at (8.10,0.92)  {once};
\node[pycap] at (12.15,0.92) {every session};

% the return path: a later session rejoins at activation, not at install
\draw[pygo, draw=MLOrange]
  (work.south) -- ++(0,-0.85) -| (act.south);
\node[pynote, text=MLOrange] at (8.10,-1.15)
  {next session: \textbf{reactivate}, do not reinstall};

\end{tikzpicture}
\end{mlfigure}

\begin{keyidea}
\textbf{Activate first, then install.} Keep course packages inside the \textbf{project environment}, not in a shared global installation.
\end{keyidea}

\begin{mlfigure}
\begin{tikzpicture}[
  font=\small,
  envbox/.style={draw=MLBlue, fill=MLPaleBlue, rounded corners=2pt,
                 align=center, inner sep=4pt, minimum width=4.3cm},
  sysbox/.style={draw=MLMuted, fill=MLPaleGray, rounded corners=2pt,
                 align=center, inner sep=4pt, minimum width=4.3cm},
  pkg/.style={font=\scriptsize\ttfamily, align=center},
  lbl/.style={font=\scriptsize\color{MLMuted}, align=center},
  hit/.style={-{Stealth[length=2.2mm]}, thick, draw=MLTeal}]

\node[sysbox] (sys) at (0,2.45) {system Python};
\node[pkg,text=MLMuted,anchor=west] at (2.35,2.45) {leave this alone};

\node[envbox] (a) at (-2.80,0) {\texttt{THIS DIRECTORY/.venv}};
\node[pkg,anchor=north] at (-2.80,-0.52) {numpy\\pandas\\scikit-learn\\jupyter};

\node[envbox] (b) at (2.80,0) {\texttt{some other project/.venv}};
\node[pkg,anchor=north] at (2.80,-0.52) {its own versions\\of the same names};

\draw[hit] (sys) -- (a);
\draw[hit] (sys) -- (b);

\node[lbl,text width=\linewidth,inner xsep=0pt] at (0,-2.60)
  {\textbf{\color{MLNavy}Separate environments, separate packages.}
   Activate the intended \textbf{project environment} before running \texttt{pip install}.};

\end{tikzpicture}
\end{mlfigure}

\section{Route B: Anaconda or Miniconda}

\textbf{Miniconda} supplies Python and \texttt{conda}; \textbf{Anaconda} adds Navigator and preinstalled packages. Download \href{https://docs.anaconda.com/miniconda/}{\textbf{Miniconda}} or \href{https://www.anaconda.com/download}{\textbf{Anaconda}}.

\begin{description}[leftmargin=1.15in,style=nextline]
  \item[Windows] Use default installer options without adding conda to PATH. Open \textbf{Anaconda Prompt} from Start.
  \item[macOS / Linux] Run the \texttt{.pkg} or \texttt{.sh} installer, then close and reopen the terminal. The prompt shows \texttt{(base)}, which means conda is active.
\end{description}

\textbf{Create} one environment for the course, with a Python version the libraries support, then \textbf{activate} it:

\begin{lstlisting}[style=mlterminal]
conda create -n ml python=3.12
conda activate ml
\end{lstlisting}

\textbf{Install dependencies} with the same \texttt{requirements.txt} used in Route A:

\begin{lstlisting}[style=mlterminal]
pip install -r requirements.txt
\end{lstlisting}

Create once, activate each session, and avoid installing into \texttt{(base)}. Check the active environment and Python:

\begin{lstlisting}[style=mlterminal]
conda env list
python -c "import sys; print(sys.executable)"
\end{lstlisting}

\begin{pitfall}
Avoid installing the same package with both conda and pip. Use pip with the course requirements. Recreate a broken environment with \texttt{conda env remove -n ml}, then repeat setup.
\end{pitfall}

\begin{troubleshoot}
\textbf{Conda not found on Windows:} use Anaconda Prompt, or initialize PowerShell there with \texttt{conda init powershell}. \textbf{Jupyter import failure:} activate \texttt{ml}, restart Jupyter, and check the kernel with Lab~1's environment cell.
\end{troubleshoot}

\section{Set up Git and GitHub for submission}

The \href{https://git-scm.com/install/}{\textbf{official Git download page}} lists current installers and package-manager commands.

\subsection{Install Git}

Use the command for your operating system, then reopen the terminal:

\begin{lstlisting}[style=mlterminal]
# Windows PowerShell
winget install --id Git.Git -e --source winget

# macOS: install Apple's command-line tools
brew install git

# Linux
sudo apt update
sudo apt install git
\end{lstlisting}

On Windows, use \textbf{Git Bash} for the SSH commands below. Confirm the installation on any system:

\begin{lstlisting}[style=mlterminal]
git --version
\end{lstlisting}

\subsection{Configure your commit identity}

Set the name and email that should appear on future commits. Use an email
connected to your GitHub account, or your GitHub-provided private email.

\begin{lstlisting}[style=mlterminal]
git config --global user.name "Your Name"
git config --global user.email "you@example.com"
git config --global init.defaultBranch main

git config --global --list
\end{lstlisting}

\subsection{Create an SSH key for GitHub}

SSH lets Git authenticate without asking for a GitHub password on every push.
First check whether this computer already has a key:

\begin{lstlisting}[style=mlterminal]
ls -al ~/.ssh
\end{lstlisting}

If \texttt{id\_ed25519} and \texttt{id\_ed25519.pub} do not exist, create them.
Enter the email connected to GitHub and choose a passphrase when prompted:

\begin{lstlisting}[style=mlterminal]
ssh-keygen -t ed25519 -C "you@example.com"
\end{lstlisting}

Press Enter to accept the default file location. Then load the \emph{private}
key into the SSH agent:

\begin{lstlisting}[style=mlterminal]
# Windows Git Bash or Linux
eval "$(ssh-agent -s)"
ssh-add ~/.ssh/id_ed25519

# macOS
eval "$(ssh-agent -s)"
ssh-add --apple-use-keychain ~/.ssh/id_ed25519
\end{lstlisting}

Copy the \emph{public} key. Never upload or share the file without
\texttt{.pub}.

\begin{lstlisting}[style=mlterminal]
# Windows Git Bash
clip < ~/.ssh/id_ed25519.pub

# macOS
pbcopy < ~/.ssh/id_ed25519.pub

# Linux: display it, then copy the complete line
cat ~/.ssh/id_ed25519.pub
\end{lstlisting}

In GitHub, open \textbf{Settings $>$ SSH and GPG keys $>$ New SSH key}. Give the
key a name such as \emph{Course laptop}, paste it, and select
\textbf{Authentication Key}. The full procedure is maintained in the
\href{https://docs.github.com/en/authentication/connecting-to-github-with-ssh}{\textbf{GitHub SSH guide}}

\subsection{Video walkthroughs for SSH setup}

Use the walkthrough for your operating system if you prefer to follow the SSH
setup on screen:

\begin{tcolorbox}[
  enhanced,
  breakable,
  colback=UTPPaleOrange,
  colframe=UTPOrange,
  colbacktitle=UTPOrange,
  boxrule=1pt,
  arc=2mm,
  title={Choose your operating system},
  fonttitle=\bfseries\color{white},
  left=3mm,right=3mm,top=2mm,bottom=2mm
]
\begin{description}[leftmargin=1.15in,style=nextline]
  \item[Linux]
    \href{https://youtu.be/AOM8CEyu3Ug?t=121}{Set up GitHub SSH on Linux}
  \item[Windows]
    \href{https://youtu.be/YT5YimhyX18?t=5}{Set up GitHub SSH on Windows}
  \item[macOS]
    \href{https://youtu.be/ACvE5Uy7ZB0?t=437}{Set up GitHub SSH on macOS}
\end{description}
\end{tcolorbox}

Do not copy another person's email address, key, username, or repository name from the video; replace each example with your own details.

\subsection{Prepare the course repository}

Do not upload the virtual environment, passwords, tokens, or API keys. In the project root, create a file named \texttt{.gitignore} containing:

\begin{lstlisting}[style=mlterminal]
.venv/
__pycache__/
*.pyc
.ipynb_checkpoints/
.env
\end{lstlisting}

Create a new repository on GitHub. Use a name such as
\texttt{utp-machine-learning}, select the visibility required by the
instructor, and leave \textbf{README}, \textbf{.gitignore}, and
\textbf{license} unchecked because the local project already exists.

From the project root, create the first commit and connect it to the empty
GitHub repository. Replace \texttt{USERNAME} and \texttt{REPOSITORY}:

\begin{lstlisting}[style=mlterminal]
git init -b main
git add .
git status
git commit -m "Add machine-learning course work"

git remote add origin git@github.com:USERNAME/REPOSITORY.git
git remote -v
git push -u origin main
\end{lstlisting}

\begin{pitfall}
Always read \texttt{git status} before committing. If a dataset, password,
token, \texttt{.env} file, or virtual environment appears, stop and update
\texttt{.gitignore} before running \texttt{git commit}. Deleting a secret in a
later commit does not remove it from earlier repository history.
\end{pitfall}

\subsection{Submit later changes}

Use short commits as each lab progresses:

\begin{lstlisting}[style=mlterminal]
git status
git add .
git commit -m "Complete Lab 1 feature preparation"
git push
\end{lstlisting}

Submit the GitHub repository address requested by the instructor and confirm that the latest commit and required notebooks are visible on GitHub.
